Fix \(\mathcal M\in\mathfrak F(\mathcal D)\) and suppose \(b\in\operatorname{col}(B_x)\).  There is \(\lambda\in\mathbb R^{\lvert\mathcal E\rvert}\) with \(B_x\lambda=b\).  Theorem~\ref{thm:target-span-iff}, Definition~\ref{def:unconditional-weight-map}, and Lemma~\ref{lem:observable-factorization} give a fixed vector \(\kappa\in\mathbb R^{\lvert\mathcal E\rvert}\) satisfying
\[
H_x\kappa=b,
\qquad
z_{x,y}^{\top}\kappa=\theta_{x,y}.
\tag{1}
\]
The division defining \(\kappa\) is valid because the latent-shift factorization and strict primitive positivity give
\[
P_O(E=e,X=x)=\pi_e\sum_{u\in\mathcal U}S_{ue}a_x(u)>0
\qquad(e\in\mathcal E).
\tag{2}
\]
No bound on \(\lVert\kappa\rVert_1\) is imposed.

We first derive the required scalar concentration inequality.  Let \(Z\) be Bernoulli with mean \(p_Z\), and put
\[
\psi(s)=\log\mathbb E\!\left[\exp\{s(Z-p_Z)\}\right].
\]
Under exponential tilting, \(\psi''(s)=p_s(1-p_s)\le1/4\), while \(\psi(0)=\psi'(0)=0\).  Two integrations yield
\[
\psi(s)\le\frac{s^2}{8}
\qquad(s\in\mathbb R).
\tag{3}
\]
For independent Bernoulli variables \(Z_1,\ldots,Z_N\) with common mean \(p_Z\), Markov's inequality and (3) imply, for \(a>0\) and \(s>0\),
\[
\begin{aligned}
\Pr\!\left\{N^{-1}\sum_{i=1}^N Z_i-p_Z\ge a\right\}
&\le \exp(-sNa)\prod_{i=1}^N
 \mathbb E\!\left[\exp\{s(Z_i-p_Z)\}\right]\\
&\le \exp\!\left(-sNa+\frac{Ns^2}{8}\right).
\end{aligned}
\tag{4}
\]
Taking \(s=4a\) and repeating the argument for the lower tail gives
\[
\Pr\!\left\{\left|N^{-1}\sum_{i=1}^N Z_i-p_Z\right|>a\right\}
\le2e^{-2Na^2}.
\tag{5}
\]

By Assumption~\ref{ass:source-iid-sampling}, every coordinate of \(\widehat H_x\) and \(\widehat z_{x,y}\) is an average of independent Bernoulli indicators with respective means \(H_x(w,e)\) and \(z_{x,y}(e)\).  By Assumption~\ref{ass:target-iid-sampling}, every coordinate of \(\widehat b\) is an average of independent Bernoulli indicators with mean \(b_w\).  Define
\[
\begin{aligned}
\mathcal A=\{&\max_{w,e}|\widehat H_x(w,e)-H_x(w,e)|\le r_s,\\
&\max_e|\widehat z_{x,y}(e)-z_{x,y}(e)|\le r_s,\\
&\max_w|\widehat b_w-b_w|\le r_t\}.
\end{aligned}
\tag{6}
\]
There are \(L=m\lvert\mathcal E\rvert+\lvert\mathcal E\rvert+m\) coordinates.  By (5) and the definitions of \(r_s,r_t\), each relevant failure probability is at most
\[
2e^{-2n_s r_s^2}=\frac{\alpha}{L}
\quad\text{or}\quad
2e^{-2n_t r_t^2}=\frac{\alpha}{L}.
\tag{7}
\]
A union bound, requiring no independence among coordinate events, gives
\[
\mathbb P_{\mathcal M}^{\,n_s,n_t}(\mathcal A)\ge1-\alpha.
\tag{8}
\]

On \(\mathcal A\), use the population vector \(\kappa\) from (1).  For every \(w\in\mathcal W\),
\[
\begin{aligned}
| (\widehat H_x\kappa-\widehat b)_w |
&\le |((\widehat H_x-H_x)\kappa)_w|+|b_w-\widehat b_w|\\
&\le r_s\sum_{e\in\mathcal E}|\kappa(e)|+r_t
=r_s\lVert\kappa\rVert_1+r_t,
\end{aligned}
\tag{9}
\]
where \(H_x\kappa=b\) was used.  Likewise,
\[
|\widehat z_{x,y}^{\top}\kappa-\theta_{x,y}|
=| (\widehat z_{x,y}-z_{x,y})^{\top}\kappa |
\le r_s\lVert\kappa\rVert_1.
\tag{10}
\]
Equations (9)--(10) are exactly the feasibility inequalities in Definition~\ref{def:concentration-projection-set} for \(\vartheta=\theta_{x,y}\).  Since \(\theta_{x,y}\in[0,1]\),
\[
\mathcal A\subseteq
\{\theta_{x,y}\in\widehat C_{x,y,1-\alpha}\}.
\tag{11}
\]
Combining (8) and (11) proves the coverage claim.

The numerical bound in (8) is independent of all cell probabilities, of the rank and nonzero singular values of \(B_x\), and of \(\lVert\kappa\rVert_1\); the realized norm is retained rather than bounded in (9)--(10).  Hence an infimum over any compatible positive laws satisfying the span condition preserves the same lower bound.  No rank estimate, matrix inverse, covariance inverse, singular-value floor, or balancing-norm bound enters the construction.  The definition permits the realized projection set to be empty off \(\mathcal A\), or to equal \([0,1]\), without affecting (11).